\section*{Part III: How Practice Becomes Skill}

\subsection*{29. Repetition Is Necessary in Many Skills, but Repetition Is Not a Theory}

A musician plays a passage 100 times. A programmer builds 5 similar applications. A sales trainee conducts 20 calls. The hours are real. The improvement is uncertain.

Repeated performance can increase speed, confidence, and familiarity. But the person may also automate errors, reinforce weak shortcuts, or stop learning once the task becomes comfortable. A course that celebrates ``100 practice problems'' without asking what each problem develops has measured volume, not skill acquisition.

A more informative unit is the \textbf{improvement cycle}: a sufficiently challenging task; an observable attempt; information about the discrepancy between current and desired performance; an explanation of the likely source; and another attempt under conditions that test the correction. The number of cycles matters, but their content matters more.

\subsection*{30. Deliberate Practice Is Targeted Correction, Not Just Hard Work}

The deliberate-practice literature emphasizes structured work aimed at improving specific aspects of performance, often with feedback and progressively more difficult tasks. Its importance has sometimes been distorted into the assertion that a fixed number of hours guarantees exceptional achievement. A meta-analysis across domains found that deliberate practice explains a meaningful but variable portion of performance differences, not everything [R9].

Imagine a learner weak at diagnosing evaluation leakage. Assigning another full-stack AI application may simply force them to repeat the same mistake within more code. A targeted exercise showing 3 deceptively similar split designs could be far more productive. After the weakness is repaired, reintegrate it into the whole application.

\textbf{Whole-task practice} develops integration. \textbf{Part-task practice} repairs bottlenecks. Expertise often requires moving between them. Neither endless isolated drills nor endless unguided projects is enough.

\subsection*{31. Practice Should Make the Invisible Decision Visible}

A beginner can sometimes produce the right output for the wrong reason. They choose temporal validation because the assignment template includes it, not because they recognized a temporal dependence. They add caching because a mentor said so, not because they diagnosed a latency pattern.

Good practice therefore elicits both \textbf{action} and \textbf{rationale}. Ask what they expect before running the code. Ask which failure they are trying to exclude. Request a short decision memo comparing plausible alternatives. When the reasoning is sound, fade the explanation requirement on routine details and move the attention to harder choices.

The goal is not bureaucratic documentation of every thought. It is the instructor's ability to distinguish knowledge from lucky execution and the learner's ability to observe their own reasoning.

\subsection*{32. Error Is Information Only When Someone Can Interpret It}

Errors can be productive because they reveal misconceptions. But error does not automatically teach. A learner who repeatedly receives ``incorrect'' without understanding why may infer that they are incapable, memorize a superficial workaround, or give up.

A useful error model has 5 stages: \textbf{expected result $\to$ actual result $\to$ discrepancy $\to$ plausible cause $\to$ discriminating next test}. The teacher's skill lies in diagnosing the \emph{reason} the attempt went wrong, not merely pointing to the wrong line.

Consider 3 students who all fail an evaluation-design question. 1 does not understand sampling; 1 understands sampling but ignores operational timing; 1 understands both but misreads the assignment. Identical grades conceal different learning needs. An excellent feedback system distinguishes them.

\subsection*{33. Feedback Is a Loop, Not a Comment}

Feedback can tell learners where they stand, but its greater value comes when it informs the next attempt. A meta-analysis of educational feedback finds a positive overall effect with substantial heterogeneity; the information conveyed is consequential [R10].

A helpful sequence is \textbf{goal $\to$ attempt $\to$ diagnosis $\to$ specific action $\to$ reattempt $\to$ check}. ``Good work'' communicates warmth, but little repair guidance. ``Your evaluation includes the same customer's later events in training and earlier events in testing; redraw the data timeline and propose a split that matches deployment'' supplies a causal diagnosis and next action.

Feedback can also be counterproductive if it overwhelms the learner, arrives after the opportunity to revise, supplies the entire solution, or frames a mistake as evidence of fixed personal inadequacy. The founder should measure not only how many comments were delivered but \textbf{whether subsequent work actually improved}.

\subsection*{34. Feedback Has at Least 4 Distinct Jobs}

\textbf{Correctness feedback} reports whether an answer or execution meets a criterion. \textbf{Diagnostic feedback} explains a relevant misconception or failure mechanism. \textbf{Strategic feedback} challenges the learner's approach and choice among alternatives. \textbf{Metacognitive feedback} helps them improve the way they monitor, test, and regulate their future work.

Each has different economics. Unit tests can provide immediate correctness information at scale. A well-crafted diagnostic rule can address common coding misconceptions. Complex system-design judgment may require a knowledgeable human who can understand context and competing trade-offs. Metacognitive coaching may be valuable when the student repeatedly misjudges readiness or avoids asking for help.

Don't spend scarce faculty time manually confirming what a deterministic test can already establish; don't trust deterministic tests to judge subtle professional reasoning they cannot observe.

\subsection*{35. Timing Matters Because Feedback Has to Meet a Teachable Moment}

Immediate feedback can stop a novice from rehearsing an error and help them establish a new procedure. Delayed feedback can sometimes protect independent thinking or create a useful opportunity to reconsider. There is no universally optimal delay across tasks.

The founder should ask: \emph{What could go wrong while waiting?} For a beginner learning basic syntax, an unresolved trivial bug can destroy the entire session. A rapid hint helps. For an advanced architecture decision, giving the recommended design immediately may erase the reasoning opportunity. A review after the learner commits to a proposal may be stronger.

The feedback service should therefore distinguish \textbf{blocking errors}, \textbf{conceptual errors}, and \textbf{judgment opportunities}. Fast resolution is not always the same as fast learning, but avoidable waiting rarely helps.

\subsection*{36. A Rubric Is a Map of Valued Decisions}

If a course grades projects mainly on visual polish, students will optimize visual polish. If every capstone passes regardless of evaluation quality, students learn that the evaluation standard is ceremonial.

A rubric makes priorities legible. For a technical system, criteria may include framing, validity of evaluation, implementation, evidence of failure testing, explanation of trade-offs, and independent defense. Rubrics should be specific enough to guide effort and reliable enough to support fair decisions, but not so mechanical that they reward checklists over judgment.

1 useful mechanism is \textbf{anchor artifacts}: examples of weak, adequate, and strong work with the reasoning for each judgment. Learners can compare their work against a visible standard; instructors can calibrate their expectations. Chapter 13 will treat assessment validity and reliability more fully.
